\section{A uniform incidence bound}\label{sec:incidence}

For a set $P$ of points and a set $\mathcal L$ of lines in
$\PP^2(\F_q)$, write
\begin{equation}\label{eq:inc-definition}
 I(P,\mathcal L)
 =\bigl|\{(p,\ell)\in P\times\mathcal L:p\in\ell\}\bigr|.
\end{equation}
Our aim is a power saving with constants independent of the field and
its characteristic. In Section~\ref{sec:linecover}, this estimate will
control intersections of a proposed line cover when no projective point
lies on many of its lines. This is the incidence input to the conclusion
that a small cover must contain almost a full pencil.

\begin{theorem}[Uniform incidence saving]\label{thm:incidence}
There exist absolute constants $\kappa\in(0,1/10)$ and $Q$ with the following
property. Suppose that $q\ge Q$ is a prime power and every proper subfield
of $\F_q$ has order at most $q^{1/3}$. If $P$ and $\mathcal L$ are sets of
projective points and projective lines over $\F_q$ satisfying
\begin{equation}\label{eq:inc-theorem-sizes}
 |P|,|\mathcal L|\le q^{1+\kappa},
\end{equation}
then
\begin{equation}\label{eq:inc-theorem-bound}
 I(P,\mathcal L)\le q^{3/2-\kappa}.
\end{equation}
\end{theorem}

The proof follows the incidence and sum--product strategy of
Bourgain, Katz, and Tao~\cite[Sections 3--4 and 6]{BKT2004}. Its geometric
step turns an incidence configuration with too many incidences into a
dense set of pairs whose differences and quotients range over small
sets. A graph argument then extracts one subset with both a small
difference set and a small quotient set. Finally, scalar coverings show
that such a subset cannot have size near the square root of the field
order under the proper-subfield hypothesis.

We develop the two algebraic ingredients first, then return to the
incidence configuration. The proofs include the simultaneous extraction
of additive and multiplicative control and the bounds on scalar
expressions that make the argument uniform over fields of nonprime
order.

\subsection{Finite-set estimates}

The first task is to pass from control of differences and quotients
along the edges of a dense bipartite graph to control on all pairs in
one subset. We begin with the sumset estimates needed to use the
resulting additive control.

For subsets of an abelian group, $X+Y$ and $X-Y$ denote the usual sumset
and difference set. For a nonnegative integer $j$, the notation $jY$
in this subsection means a sum of $j$ copies of $Y$, with $0Y=\{0\}$.
Later, a field element multiplying a set will always mean field scaling;
we will specify explicitly when an integer denotes an iterated sumset.

The first two estimates below are Ruzsa's triangle inequality and
covering lemma; the maximal-translates proof of the latter appears in
\cite[p.~324]{Ruzsa1999}. The multiple-sum bound is a Pl\"unnecke--Ruzsa
estimate, for which we use Petridis's minimizing-subset
proof~\cite[Proposition 2.1 and Section 3]{Petridis2012}.

\begin{lemma}[Triangle, covering, and multiple-sum estimates]
\label{inc-sumset-estimates}
Let the sets below be finite subsets of an abelian group.
\begin{enumerate}
\item If $Z$ is nonempty, then
\begin{equation}\label{eq:inc-triangle}
 |X-Y|\,|Z|\le |X-Z|\,|Z-Y|.
\end{equation}
\item If $Y$ is nonempty and $|X-Y|\le M|Y|$, then $X$ is contained in
the union of at most $M$ translates of $Y-Y$.
\item If $Y$ is nonempty and $|Y-Y|\le M|Y|$, then, for all nonnegative
integers $u,v$,
\begin{equation}\label{eq:inc-multiple-sums}
 |uY-vY|\le M^{u+v}|Y|.
\end{equation}
\end{enumerate}
\end{lemma}

\begin{proof}
For the triangle inequality, choose one representation $d=x_d-y_d$
for each $d\in X-Y$. The map
\[
 (d,z)\longmapsto (x_d-z,z-y_d)
\]
from $(X-Y)\times Z$ into $(X-Z)\times(Z-Y)$ is injective: summing the
two coordinates recovers $d$, after which the chosen representation
recovers $z$. This proves~\eqref{eq:inc-triangle}.

For the covering statement, take a maximal set $T\subseteq X$ for
which the translates $t-Y$, $t\in T$, are pairwise disjoint. Their union
is contained in $X-Y$, so $|T|\le M$. Maximality implies that for every
$x\in X$ some $x-Y$ intersects $t-Y$ with $t\in T$, whence
$x\in t+(Y-Y)$.

For the multiple-sum bound, the minimizing-subset argument of
Petridis~\cite{Petridis2012} lets us iterate a one-step growth estimate.
Put $Q_0=-Y$ and choose a nonempty
$W\subseteq Y$ minimizing
\[
 \mu=\frac{|W+Q_0|}{|W|}.
\]
Then $\mu\le M$, and every $W'\subseteq W$, including the empty set,
satisfies $|W'+Q_0|\ge\mu|W'|$. We claim that for every finite set $E_0$,
\begin{equation}\label{eq:inc-petridis-growth}
 |W+Q_0+E_0|\le\mu|W+E_0|.
\end{equation}
The assertion is immediate if $E_0$ is empty. Otherwise order its
elements as $e_1,\ldots,e_k$. When the translate $W+e_i$ is added to
the preceding translates, let
\[
 W_i'=\{w\in W:w+e_i\in\bigcup_{j<i}(W+e_j)\}.
\]
It contributes exactly $|W|-|W_i'|$ new elements. Moreover,
$W_i'+Q_0+e_i$ is contained in the union of the preceding enlarged
translates $W+Q_0+e_j$. Thus the contribution of $W+Q_0+e_i$ is at most
\[
 |W+Q_0|-|W_i'+Q_0|
 \le\mu\bigl(|W|-|W_i'|\bigr).
\]
Summing proves~\eqref{eq:inc-petridis-growth}. Iterating this estimate
gives $|W+jQ_0|\le\mu^j|W|$ for every nonnegative integer $j$.
We now remove the auxiliary set $W$ by applying
\eqref{eq:inc-triangle} to $uQ_0,vQ_0$ with intermediate set $-W$:
\[
 |uQ_0-vQ_0|\,|W|
 \le |W+uQ_0|\,|W+vQ_0|
 \le\mu^{u+v}|W|^2.
\]
Negation preserves cardinality, so $|uQ_0-vQ_0|=|uY-vY|$. Since
$|W|\le|Y|$ and $\mu\le M$, this is~\eqref{eq:inc-multiple-sums}.
\end{proof}

We next carry out the extraction from a dense graph. The argument uses
the common-neighborhood method in Gowers's quantitative refinement
\cite[Section 4, Lemma 11 and Proposition 12]{Gowers1998} of the
Balog--Szemer\'edi theorem~\cite{BalogSzemeredi1994}; see also
\cite[Section 5, Lemmas 5.1--5.2]{FoxSudakov2011}. The point is to use the same walks
for both additive and multiplicative edge labels. We can therefore
obtain both bounds on one subset, which is what the subsequent
subfield argument requires. In the geometric application, the edge
labels will come from the four pencils through four chosen points.

\begin{lemma}[Simultaneous difference and quotient extraction]
\label{inc-simultaneous-extraction}
Let $F$ be a field, let $A,B\subseteq F^*$ be finite, and let
$G\subseteq A\times B$. Define
\[
 D=\{a-b:(a,b)\in G\},\qquad
 J=\{a/b:(a,b)\in G\}.
\]
Suppose that, for real numbers $n>0$ and $V\ge2$,
\begin{equation}\label{eq:inc-graph-hypotheses}
 |A|,|B|,|D|,|J|\le Vn,
 \qquad |G|\ge n^2/V.
\end{equation}
Then there is $S\subseteq A$ such that
\begin{equation}\label{eq:inc-extracted-size}
 \frac{3n}{8V^2}\le |S|\le Vn
\end{equation}
and
\begin{equation}\label{eq:inc-extracted-expansion}
 |S-S|,|S/S|\le 2^{20}V^{22}n.
\end{equation}
Here $S/S=\{a/a':a,a'\in S\}$.
\end{lemma}

\begin{proof}
We first find a large set of vertices such that every ordered pair
is joined by many walks of length four. Counting the possible edge
labels of these walks will then bound both differences and quotients
of the endpoints.

Regard $G$ as a bipartite graph, and put
$\tau=(512V^8)^{-1}$. Call an ordered pair $(a,a_2)\in A^2$ bad if
it has fewer than $\tau n$ common neighbors in $B$. To find a
neighborhood with few bad pairs, first note that vertices of $B$
with fewer than $n/(2V^2)$ neighbors account for at most
$|B|n/(2V^2)\le n^2/(2V)$ edges. Since every vertex of $B$ has at
most $Vn$ neighbors, the set of vertices with at least $n/(2V^2)$
neighbors has size at least $n/(2V^2)$.

Summed over all $b\in B$, the number of bad ordered pairs in the
neighborhood of $b$ is at most
\[
 |A|^2\tau n\le V^2\tau n^3.
\]
Averaging over the large neighborhoods just identified, one obtains
a neighborhood $A_0$ of size at least $n/(2V^2)$ in which the number
of bad ordered pairs is at most
\[
 2V^4\tau n^2=\frac{n^2}{256V^4}
 \le \frac{|A_0|^2}{64}
 \le \frac{|A_0|^2}{16}.
\]
Let $S$ consist of those $a\in A_0$ having at most $|A_0|/4$ bad
partners in $A_0$. At most $|A_0|/4$ elements are discarded, so
$|S|\ge3|A_0|/4$, proving~\eqref{eq:inc-extracted-size}.

The choice of $S$ ensures that, for every $a,a'\in S$, at least
$|A_0|/2$ elements $a_2\in A_0$ form nonbad pairs with both endpoints.
Each such $a_2$ gives at least $(\tau n)^2$
ordered walks
\[
 a,b_1,a_2,b_2,a'
\]
whose four edges belong to $G$. Repeated vertices are allowed.
Consequently each ordered endpoint pair in $S^2$ has at least
\begin{equation}\label{eq:inc-walk-count}
 \frac{|A_0|}{2}(\tau n)^2
 \ge\frac{n^3}{2^{20}V^{18}}
\end{equation}
such walks. This lower bound now converts the limited number of edge
labels into a bound on endpoint differences.

For one of these walks, write its edge differences, always taking
the $A$ coordinate minus the $B$ coordinate, as
\[
 d_1=a-b_1,\quad d_2=a_2-b_1,\quad
 d_3=a_2-b_2,\quad d_4=a'-b_2.
\]
They satisfy $a-a'=d_1-d_2+d_3-d_4$. For a fixed endpoint pair,
the four differences determine the walk uniquely: starting from $a$,
one recovers successively $b_1,a_2,b_2,a'$. Choose one ordered
endpoint pair for each element of $S-S$. The resulting collections
of difference quadruples are disjoint for different differences,
because their alternating sums differ. Counting in $D^4$ and using
\eqref{eq:inc-walk-count} gives
\[
 |S-S|\le
 \frac{|D|^4}{n^3/(2^{20}V^{18})}
 \le 2^{20}V^{22}n.
\]

The same walks give the quotient bound on the same set $S$. Put
\[
 j_1=a/b_1,\quad j_2=a_2/b_1,\quad
 j_3=a_2/b_2,\quad j_4=a'/b_2.
\]
All are nonzero, and $a/a'=j_1j_3/(j_2j_4)$. From a fixed initial
$a$ and these four quotients one recovers
$b_1=a/j_1$, $a_2=j_2b_1$, $b_2=a_2/j_3$, and $a'=j_4b_2$.
Choosing one endpoint pair for each element of $S/S$ and counting
in $J^4$ therefore proves the other bound in
\eqref{eq:inc-extracted-expansion}.
\end{proof}

\subsection{Bounded scalar expressions and subfields}

We now turn to the algebraic obstruction: a set near the square-root
field scale cannot have both the small difference set and the small
quotient set supplied by the extraction lemma. The first step is to
bound sums involving scalars formed using a bounded number of field
operations. We use the scalar-covering method of Bourgain, Katz, and
Tao~\cite[Lemma 3.2 and Proposition 3.3]{BKT2004}, keeping the dependence
on the number of operations explicit.

An expression in a set $E\subseteq F$ is built from members of
$E\cup\{0,1\}$ using binary addition, subtraction, and multiplication,
and unary negation and inversion. Each of these operations counts
as one operation. Inversion is allowed only when its argument is
nonzero.

\begin{lemma}[Uniform scalar coverings]\label{inc-scalar-coverings}
Let $B$ be a nonempty finite subset of a field $F$, let $E\subseteq B$,
and let $K\ge2$. Suppose that
\begin{equation}\label{eq:inc-scalar-hypotheses}
 |B-B|\le K|B|,
 \qquad |B-eB|\le 2K^3|B|\quad(e\in E).
\end{equation}
Define
\begin{equation}\label{eq:inc-scalar-exponents}
 C_0=4,\qquad C_{m+1}=3C_m+5,
 \qquad\text{so that }C_m=\frac{13\cdot3^m-5}{2}.
\end{equation}
If $\lambda\in F$ has an expression in $E$ using at most $m$
operations, then $\lambda B$ is contained in at most $K^{C_m}$
translates of $B-B$. In particular,
\begin{equation}\label{eq:inc-scalar-sum-bound}
 |E+\lambda E|\le |B+\lambda B|
 \le K^{C_m+3}|B|.
\end{equation}
These bounds are uniform over all such expressions and all choices
of their entries from $E$.
\end{lemma}

\begin{proof}
Put $H=B-B$, so $H=-H$ and $0\in H$. In the remainder of this proof,
integer multiples of $B$ or $H$ inside sumsets denote iterated sums.
We prove the covering assertion by induction on expression length.
To combine coverings at an induction step, we will need a fixed
covering of $2H$ by translates of $H$. By
Lemma~\ref{inc-sumset-estimates},
\[
 |2H-B|=|2B-3B|\le K^5|B|.
\]
The covering statement of that lemma consequently supplies a set
$T_*\subseteq F$ with $|T_*|\le K^5$ such that
\begin{equation}\label{eq:inc-two-differences-cover}
 2H\subseteq T_*+H.
\end{equation}

Every $e\in E$ has a covering $eB\subseteq T_e+H$ with
$|T_e|\le2K^3\le K^4$, by
\eqref{eq:inc-scalar-hypotheses} and the covering lemma.
The scalars $0$ and $1$ also have coverings by one translate of $H$:
use $\{0\}\subseteq H$ and $B\subseteq b_0+H$ for any $b_0\in B$.
This proves the case $m=0$.

We now check how the number of translates changes under each allowed
operation. Suppose that
\[
 \lambda B\subseteq T+H,\qquad
 \mu B\subseteq U+H,\qquad
 |T|\le K^a,\quad |U|\le K^b.
\]
Negation replaces $T$ by $-T$. For addition and subtraction,
\[
 (\lambda\pm\mu)B
 \subseteq \lambda B\pm\mu B
 \subseteq T\pm U+2H
 \subseteq T\pm U+T_*+H,
\]
so the number of translates is at most $K^{a+b+5}$. For a product,
\[
 \lambda\mu B
 \subseteq \lambda U+\lambda H
 =\lambda U+(\lambda B-\lambda B)
 \subseteq\lambda U+T-T+T_*+H.
\]
This uses at most $K^{b+2a+5}$ translates and is valid even if
$\lambda=0$. Finally, if $\lambda\ne0$, then
\[
 |\lambda B-B|
 \le |T|\,|H-B|
 =|T|\,|B-2B|
 \le K^{a+3}|B|.
\]
Scaling by $\lambda^{-1}$ and then negating gives
$|\lambda^{-1}B-B|\le K^{a+3}|B|$. The covering lemma now covers
$\lambda^{-1}B$ by at most $K^{a+3}$ translates of $H$.

These bounds give the claimed uniform induction: in an expression
with at most $m+1$ operations, each immediate subexpression has at
most $m$ operations, so every resulting covering uses at most
$K^{3C_m+5}$ translates. This proves the covering assertion
with~\eqref{eq:inc-scalar-exponents}.

To deduce the sumset bound from the covering, let
$\lambda B\subseteq T+H$ be the covering just obtained. Then
\[
 |B+\lambda B|
 \le |T|\,|B+H|
 =|T|\,|2B-B|
 \le K^{C_m+3}|B|,
\]
again by Lemma~\ref{inc-sumset-estimates}. Since $E\subseteq B$,
this also bounds $|E+\lambda E|$.
\end{proof}

When working along a sequence $q\to\infty$, a positive quantity
$K=q^{o(1)}$ means that $\log K/\log q\to0$. Likewise,
$|S|=q^{1/2+o(1)}$ means that $\log|S|/\log q\to1/2$.
In particular, every fixed power of a subpower factor is still a
subpower factor. This is why a uniform bound for expressions of
fixed length suffices for the next lemma.

\begin{lemma}[Exclusion of an intermediate approximate subfield]
\label{inc-no-approximate-subfield}
There is no sequence of prime powers $q\to\infty$, each satisfying
the proper-subfield hypothesis of Theorem~\ref{thm:incidence}, together
with sets $S\subseteq\F_q^*$ and numbers $K\ge2$ such that
\begin{equation}\label{eq:inc-approximate-subfield-hypotheses}
 |S|=q^{1/2+o(1)},\qquad K=q^{o(1)},\qquad
 |S-S|,|S/S|\le K|S|.
\end{equation}
\end{lemma}

\begin{proof}
Suppose that such a sequence exists. We first use multiplicative
overlaps to obtain sets satisfying the scalar-covering hypotheses.
Ratios of differences in a large subset will then form a subfield.
The proper-subfield hypothesis will force that subfield to be the
whole field, where averaging will contradict the scalar-covering
bound.

\paragraph{Extracting scalar-covering hypotheses.}
Put $s=|S|$. The $s$ sets
$S/a$, $a\in S$, each have size $s$, and their union is contained
in $S/S$. Cauchy--Schwarz gives
\begin{equation}\label{eq:inc-multiplicative-overlaps}
 \sum_{a,b\in S}|S/a\cap S/b|
 \ge\frac{s^4}{|S/S|}\ge\frac{s^3}{K}.
\end{equation}
Thus some $b\in S$ satisfies
$\sum_{a\in S}|S/a\cap S/b|\ge s^2/K$.
Since every summand is at most $s$, at least $s/(2K)$ values
$a\in S$ have $|S/a\cap S/b|\ge s/(2K)$: otherwise the sum of
the large summands would be less than $s^2/(2K)$ and that of all
the other summands at most $s^2/(2K)$.

Let $E$ be the set of $e=a/b$ for these values of $a$, and put
$B=S/b$. Then
\begin{equation}\label{eq:inc-overlap-subset}
 E\subseteq B,\qquad |B|=s,\qquad
 |E|\ge\frac{s}{2K},\qquad
 |S\cap eS|\ge\frac{s}{2K}\quad(e\in E).
\end{equation}
The last inequality follows by multiplying $S/a\cap S/b$ by $a$.
For $Z=S\cap eS$, the triangle inequality and the containments
$Z\subseteq S,eS$ imply
\[
 |S-eS|
 \le\frac{|S-Z|\,|Z-eS|}{|Z|}
 \le\frac{(Ks)^2}{s/(2K)}=2K^3s.
\]
After scaling by $1/b$, we have
\[
 |B-B|\le Ks,\qquad |B-eB|\le2K^3s\quad(e\in E).
\]
Thus Lemma~\ref{inc-scalar-coverings} applies to $B$ and $E$.
We have obtained uniform control for every scalar expressible using
a bounded number of operations on elements of $E$.

\paragraph{Ratios of differences form a subfield.}
Write $t=|E|$. By~\eqref{eq:inc-overlap-subset} and
\eqref{eq:inc-approximate-subfield-hypotheses},
$t=q^{1/2+o(1)}$. In particular, $t\ge2$ for all sufficiently
large fields in the sequence. Define
\begin{equation}\label{eq:inc-ratio-set}
 R_0=\left\{\frac{x-x'}{y-y'}:
 x,x',y,y'\in E,\ y\ne y'\right\}.
\end{equation}
To prove closure under addition and multiplication, we apply the
scalar-covering bound to the result of either operation.
Every member of $R_0$ has an expression in $E$ with at most four
operations. Consequently, for $r,r'\in R_0$, each of $r+r'$ and
$rr'$ has an expression with at most nine operations. For either
candidate $\lambda$, Lemma~\ref{inc-scalar-coverings} gives
\[
 |E+\lambda E|\le K^{C_9+3}s<t^2
\]
for all sufficiently large $q$, uniformly over $r,r'$.
Indeed, the ratio of this upper bound to $t^2$ is at most
$4K^{C_9+5}/s$, which tends to zero. Therefore the map
$(x,y)\mapsto x+\lambda y$ from $E^2$ is not injective. A collision
between distinct pairs satisfies $y\ne y'$, since otherwise also
$x=x'$, and gives
\[
 \lambda=\frac{x'-x}{y-y'}\in R_0.
\]
It follows that $R_0$ is closed under addition and multiplication.
It contains $0$ and $1$ and, directly from its definition, is closed
under negation and reciprocals of nonzero elements. Thus $R_0$ is a
subfield of $\F_q$.

The size of $E$ now prevents this subfield from being proper.
Choose distinct $e_0,e_1\in E$. The map
\[
 e\longmapsto\frac{e-e_0}{e_1-e_0}
\]
injects $E$ into $R_0$, so $|R_0|\ge t>q^{1/3}$ for large $q$.
The proper-subfield hypothesis forces $R_0=\F_q$.

\paragraph{A scalar with a large sumset.}
Every field element consequently has a four-operation representation
using elements of $E$. The scalar-covering estimate therefore bounds
$|E+\lambda E|$ for every scalar. To reach a contradiction, we average
over the field and find a scalar for which this sumset is too large.

For $\lambda\in\F_q$, let $N_\lambda$ be the number of ordered
quadruples $(x,y,x',y')\in E^4$ satisfying
$x+\lambda y=x'+\lambda y'$. There are $t^2$ identical pairs,
which give a collision for every $\lambda$. If $y\ne y'$, exactly
one value of $\lambda$ gives a collision; if $y=y'$ but $x\ne x'$,
none does. Hence
\begin{equation}\label{eq:inc-average-collisions}
 \frac1q\sum_{\lambda\in\F_q}N_\lambda
 =t^2+\frac{t^4-t^3}{q}
 \le t^2+\frac{t^4}{q}.
\end{equation}
For some $\lambda$, Cauchy--Schwarz applied to the fibers of
$(x,y)\mapsto x+\lambda y$ therefore yields
\begin{equation}\label{eq:inc-large-dilate-sum}
 |E+\lambda E|
 \ge\frac{t^4}{t^2+t^4/q}
 \ge\frac12\min(t^2,q)
 =q^{1-o(1)}.
\end{equation}
But $R_0=\F_q$, so this particular $\lambda$ itself has a
four-operation representation in~\eqref{eq:inc-ratio-set}.
Lemma~\ref{inc-scalar-coverings} bounds its sumset by
$K^{C_4+3}s=q^{1/2+o(1)}$, contradicting
\eqref{eq:inc-large-dilate-sum}. The expression length here remains
bounded; no iteration of field generation is being used.
\end{proof}

\Needspace{10\baselineskip}
\subsection{Proof of the incidence theorem}

The algebraic ingredients are now in place. It remains to obtain the
hypotheses of Lemma~\ref{inc-simultaneous-extraction} from a configuration
with too many incidences. We use the projective-coordinate reduction
in the incidence strategy of Bourgain, Katz, and
Tao~\cite[Section 6]{BKT2004}: after controlling degrees, we choose four
points whose pencils encode the two coordinates, their difference,
and their quotient.

\begin{proof}[Proof of Theorem~\ref{thm:incidence}]
We argue by contradiction, allowing all quantitative losses to be
subpower factors along a sequence of fields. Suppose the theorem is
false. For a sequence of positive exponents
$\kappa_j\to0$ in $(0,1/10)$, choose admissible counterexamples with
$q_j\to\infty$. Such choices are possible because the failure of
the theorem says that for each fixed exponent there are arbitrarily
large counterexamples. Along this sequence, put
$U=\max(2,q_j^{\kappa_j})$ and suppress the index $j$. We have
\begin{equation}\label{eq:inc-counterexample-sequence}
 |P|,|\mathcal L|\le Uq,\qquad
 I(P,\mathcal L)\ge\frac{q^{3/2}}{U},\qquad
 U\ge2,\qquad U=q^{o(1)}.
\end{equation}
All bounds below use absolute constants. Every assertion made for
sufficiently large $q$ is along this sequence. Our goal is to construct
from these counterexamples sets of the size and small-expansion form
excluded by Lemma~\ref{inc-no-approximate-subfield}.

\paragraph{Pruning the incidence graph.}
We first retain many incidences while putting both an upper and a
lower bound on every remaining degree. Consider the bipartite graph
on $P$ and $\mathcal L$ whose edges are incidences. If $d(v)$ denotes a degree on either side, then counting
ordered pairs on the opposite side and using the unique join of two
distinct projective points, or the unique intersection of two
distinct projective lines, gives
\begin{equation}\label{eq:inc-degree-squares}
 \sum_v d(v)^2\le2U^2q^2
 \quad\text{on each side}.
\end{equation}
Here $\sum_vd(v)(d(v)-1)\le U^2q^2$, and
$\sum_vd(v)=I(P,\mathcal L)\le U^2q^2$.

Set
\begin{equation}\label{eq:inc-degree-cutoffs}
 H=16U^3q^{1/2},\qquad d=\frac{q^{1/2}}{8U^2}.
\end{equation}
Delete every vertex whose initial degree exceeds $H$. By
\eqref{eq:inc-degree-squares}, the total number of incidences lost,
counting both sides, is at most
\[
 \frac{4U^2q^2}{H}=\frac{q^{3/2}}{4U}.
\]
Next, iteratively delete vertices whose current degree is less than
$d$. Each of the at most $2Uq$ original vertices is deleted at most
once, so this stage loses at most
$2Uq\,d=q^{3/2}/(4U)$ incidences. Denote the surviving sets again by
$P,\mathcal L$. They satisfy
\begin{equation}\label{eq:inc-pruned-bounds}
 I(P,\mathcal L)\ge\frac{q^{3/2}}{2U},\qquad
 \frac{q}{32U^4}\le |P|,|\mathcal L|\le Uq,
 \qquad d\le d(v)\le H.
\end{equation}
The size lower bounds follow by dividing the surviving incidence
count by the maximum degree $H$.

The pruning also controls the intersection of $P$ with
\emph{every} projective line, including lines not in $\mathcal L$.
We will need this stronger form of control when discarding points
on a few lines chosen later. Fix a projective line $\ell$. Each point
of $P\cap\ell$ is incident with at least $d-1$ lines of $\mathcal L$ different from $\ell$.
These lines are distinct for different points of $P\cap\ell$.
Since $d\to\infty$, for large $q$ we obtain
\begin{equation}\label{eq:inc-all-line-cap}
 |P\cap\ell|
 \le\frac{|\mathcal L|}{d-1}
 \le16U^3q^{1/2}=H
 \qquad\text{for every projective line }\ell.
\end{equation}

\paragraph{Four anchors with many common witnesses.}
We seek four points, three collinear and the fourth off their line,
together with many points joined to each of them by lines of
$\mathcal L$. We obtain the four points by first counting such
configurations around each possible common point and then averaging.

For $x\in P$, let $N(x)$ consist of the points in $P\setminus\{x\}$
joined to $x$ by a line of $\mathcal L$. Distinct lines through $x$
have disjoint point sets away from $x$, so
\begin{equation}\label{eq:inc-neighborhood-size}
 |N(x)|\ge d(d-1)\ge\frac{q}{128U^4}
\end{equation}
for large $q$. Put $t_\ell=|N(x)\cap\ell|$. Summing degrees of the
points of $N(x)$ gives
\[
 \sum_{\ell\in\mathcal L}t_\ell
 \ge d|N(x)|\ge\frac{q^{3/2}}{1024U^6}.
\]
Lines with $t_\ell\le2$ contribute at most $2Uq$. Because $U=q^{o(1)}$,
for sufficiently large $q$ this is at most half of the displayed
lower bound. Consequently
\[
 \sum_{\ell:t_\ell\ge3}t_\ell
 \ge\frac{q^{3/2}}{2048U^6}.
\]
For every integer $t\ge3$, one has
$t(t-1)(t-2)\ge(2/9)t^3$. H\"older's inequality therefore shows
that the number of ordered triples of distinct points of $N(x)$
lying on a line of $\mathcal L$ is at least
\begin{equation}\label{eq:inc-collinear-triples}
 \begin{split}
 \sum_{\ell\in\mathcal L}t_\ell(t_\ell-1)(t_\ell-2)
 &\ge\frac{2}{9|\mathcal L|^2}
       \left(\sum_{\ell:t_\ell\ge3}t_\ell\right)^3\\
 &\ge\frac{q^{5/2}}{9\cdot2^{32}U^{20}}.
 \end{split}
\end{equation}
By~\eqref{eq:inc-all-line-cap} and
\eqref{eq:inc-neighborhood-size}, each such triple admits at least
\[
 \frac{q}{128U^4}-H\ge\frac{q}{256U^4}
\]
choices of a fourth point in $N(x)$ outside its line, once $q$ is
large enough. Summing over $x$ and using
\eqref{eq:inc-pruned-bounds}, we obtain at least
\begin{equation}\label{eq:inc-witness-quadruple-count}
 \frac{q^{9/2}}{9\cdot2^{45}U^{28}}
\end{equation}
ordered choices $(x,p_1,p_2,p_3,p_4)$ with $p_1,p_2,p_3$ distinct
and collinear on a line of $\mathcal L$, with $p_4$ outside that
line, and with all four $p_i$ in $N(x)$.

We now reverse the count to fix the four anchor points while retaining
many choices of the common point $x$. The number of possible ordered
quadruples $(p_1,p_2,p_3,p_4)$ in this count is at most
\begin{equation}\label{eq:inc-quadruple-upper-bound}
 |\mathcal L|H^3|P|
 \le2^{12}U^{11}q^{7/2}.
\end{equation}
Some quadruple thus has at least
\begin{equation}\label{eq:inc-common-witnesses}
 c_0\frac{q}{U^{39}}
 \quad\text{witnesses }x,\qquad
 c_0=\frac1{9\cdot2^{57}}.
\end{equation}
Discard the witnesses on the common line of $p_1,p_2,p_3$ and on
the joins of $p_1,p_4$ and of $p_2,p_4$. The cap
\eqref{eq:inc-all-line-cap} discards at most $3H$ witnesses.
Since $U=q^{o(1)}$, this is eventually less than
$c_0q/(2U^{39})$. At least that many witnesses remain. Their joins
to the four anchors will supply the graph and the four small sets
required by Lemma~\ref{inc-simultaneous-extraction}.

\paragraph{The four pencils as coordinates.}
Choose projective coordinates in which
\[
 p_1=(1:0:0),\qquad p_2=(0:1:0),\qquad
 p_3=(1:1:0),\qquad p_4=(0:0:1).
\]
This is possible over $\F_q$: representatives of $p_1,p_2$ span
their line, the representative of the distinct third point has
both coefficients nonzero and permits rescaling of the first two,
and a representative of $p_4$ completes a basis. The discarded
lines are precisely the coordinate line at infinity and the two
coordinate axes through $p_4$. Hence every remaining witness has
coordinates $(a:b:1)$ with $a,b\ne0$.

Let $G$ be the set of these pairs $(a,b)$, and let $A,B$ be its
coordinate projections. Define $D$ and $J$ from $G$ as in
Lemma~\ref{inc-simultaneous-extraction}. The chosen coordinates make
the roles of the four pencils explicit: for a witness $(a:b:1)$,
the joins to $p_2,p_1,p_3,p_4$, respectively, have equations
\[
 X=aZ,\qquad Y=bZ,\qquad
 X-Y=(a-b)Z,\qquad X=(a/b)Y.
\]
Each join is a member of $\mathcal L$, by the witness property.
Each anchor belongs to $P$ and has degree at most $H$.
Distinct values of the corresponding parameter give distinct
lines in its pencil. Thus
\[
 |A|,|B|,|D|,|J|\le H,
 \qquad |G|\ge\frac{c_0q}{2U^{39}}.
\]
With
\begin{equation}\label{eq:inc-graph-parameters}
 n=q^{1/2},\qquad V=\frac{2}{c_0}U^{39},
\end{equation}
we have $V\ge2$, $V=q^{o(1)}$, and all the hypotheses
\eqref{eq:inc-graph-hypotheses} hold.

The incidence configuration now meets the extraction hypotheses.
Lemma~\ref{inc-simultaneous-extraction} gives
$S\subseteq\F_q^*$ satisfying
\[
 \frac{3q^{1/2}}{8V^2}\le |S|\le Vq^{1/2},\qquad
 |S-S|,|S/S|\le2^{20}V^{22}q^{1/2}.
\]
In particular $|S|=q^{1/2+o(1)}$, and, on setting
$K=2^{23}V^{24}$, we have $K\ge2$, $K=q^{o(1)}$, and
$|S-S|,|S/S|\le K|S|$. This contradicts
Lemma~\ref{inc-no-approximate-subfield}.

We have excluded every sequence of counterexamples of the form
\eqref{eq:inc-counterexample-sequence}. The negation of the theorem
would have produced such a sequence, so some absolute
$\kappa\in(0,1/10)$ and $Q$ satisfy the stated conclusion.
\end{proof}
